\section{Communication and Computation Cost}
\label{sec:cost}

We compare \Cref{cons:bitop} against the naive arithmetization that
materialises a fresh trace column $w := \bitop_*(v)$, commits to it,
and proves $w_i = \bitop(v_i)$ for every row $i$.

\subsection{Per-spec accounting}

Per declared spec $(j_\ell,\bitop_\ell) \in \Sbit$:

\begin{table}[h]
\centering
\renewcommand{\arraystretch}{1.2}
\begin{tabular}{l c c c c}
\toprule
& Commitment & Sumcheck rounds & Field elements (proof) & New challenges \\
\midrule
Naive (commit-and-equate)
  & $1$ column
  & $\Theta(\log n)$ extra
  & $\Theta(\log n)$ + opening
  & $\ge 1$ \\
This construction
  & $0$
  & $0$
  & $1$
  & $0$ \\
\bottomrule
\end{tabular}
\caption{Per-spec cost. The naive count assumes the equality $w_i =
\bitop(v_i)$ is proven by a separate sumcheck or batched into the
existing aggregation.}
\end{table}

\subsection{Verifier-side computation}

Verifier-side, per spec, the construction adds:
\begin{itemize}[leftmargin=2em]
  \item One $\F_q$-linear permutation of $32$ coefficients
    ($\bitop_\ell$ on $\liftring$): $O(32)$ assignments;
  \item One $\psia$-projection of a degree-$<32$ polynomial in
    $\liftring$: $32$ multiplications and $31$ additions in $\F_q$
    (or $O(32)$ multiplications if $\alpha$-powers are precomputed);
  \item One substitution into the mp\_eval subclaim equation: $O(1)$
    field operations.
\end{itemize}
Total: $O(32)$ $\F_q$-operations per spec, independent of $n$.

\subsection{Prover-side computation}

Prover-side, per spec, the construction adds:
\begin{itemize}[leftmargin=2em]
  \item Materialisation of $M_\ell$: $n$ rows, each $O(32)$
    $\F_q$-operations to apply $\psia \circ \bitop_\ell$;
  \item Inclusion of $M_\ell$ in the constraint-aggregation sumcheck:
    one extra MLE in the sumcheck working set, contributing $O(n)$
    additional folding work per round (i.e., $O(n \nu)$ total over
    the $\nu$ rounds), comparable to one extra row-shift virtual
    column;
  \item Inclusion of $M_\ell$ as a source in mp\_eval: one extra
    source claim, contributing the standard mp\_eval per-source
    overhead.
\end{itemize}

\subsection{Asymptotic comparison}

For an arithmetization with $\sigma$ instances of bit-ops over a
common source column (e.g., SHA-256 with three instances each in
$\sigma_0$ and $\sigma_1$, repeated across $48$ message-schedule
rounds), the construction reduces:
\begin{itemize}[leftmargin=2em]
  \item committed column count by $\Theta(\sigma)$;
  \item proof-byte overhead from $\Theta(\sigma \log n)$ to
    $\Theta(\sigma)$ field elements;
  \item the equality-constraint count by exactly $\sigma$.
\end{itemize}
Because the witness commitment dominates honest-prover proof size in
typical instantiations of the host PIOP, this cost reduction is
proportional to the bit-op density of the arithmetization, not to its
total constraint count.

\subsection{Concrete numbers in the implementation}

In the deployed Zinc+ SHA-256 UAIR (see the SHA UAIR
specification), the use of the construction yields:
\[
  m \;=\; 6, \qquad
  |\Sup\!\restriction_{\mathrm{binary}}|\;=\;19,
\]
i.e., the SHA UAIR declares six bit-op virtual columns over a
$19$-column binary-poly section of the trace. The bit-op virtual
columns add exactly $6$ field elements to the proof and $0$ rounds.
